\documentclass[10pt,a4paper]{amsart}
\usepackage[margin=19mm]{geometry}
\usepackage{amsmath,amssymb,mathtools}
\usepackage[scr=rsfs,cal=euler]{mathalfa}
\usepackage{xfrac}
\usepackage[unicode,hidelinks,bookmarks=false]{hyperref}
\newcommand{\ie}{i.e.\ }
\newcommand{\cf}{cf.\ }
\newcommand{\resp}{respectively\ }
\newcommand{\Subsection}{\subsection}
\providecommand{\nobreakdash}{\nobreak\hbox{-}}
\setlength{\emergencystretch}{2em}
\multlinegap0pt
\usepackage{xcolor}
\usepackage{amssymb}
\usepackage{enumerate}
\usepackage{mathtools}
\usepackage[shortlabels]{enumitem}
\newtheorem{lemma}{Lemma}
\newcommand{\R}{\mathbb{R}}
\title{Martingale Transforms and Compensated Bellman Estimates for Dunkl Riesz Transforms}
\author{Francesco D'Emilio, Brett D. Wick}
\date{Source: arXiv:2609.12117v1 (CC BY 4.0)}
\begin{document}
\maketitle

\noindent\emph{Source and licence.} Martingale Transforms and Compensated Bellman Estimates for Dunkl Riesz Transforms. Version arXiv:2609.12117v1 (CC BY 4.0), distributed by arXiv under the Creative Commons Attribution 4.0 International licence, http://creativecommons.org/licenses/by/4.0/, as recorded in the arXiv licence metadata for this exact version and shown on the abstract page https://arxiv.org/abs/2609.12117v1. Full licence notice: \texttt{licenses/arxiv-2609.12117-CC-BY-4.0.txt}.\par\smallskip

\noindent\textit{Editorial scope.} Counted unit: the article's lemma lemma (source0.tex lines 2679--2685). The article's complete original proof of it is delivered with it (source0.tex lines 2687--2873, one \texttt{proof} environment of 187 lines). No mathematics is written, repaired or re-derived: the statement and the proof are the article's own lines, in the article's own notation, and no retained line is substituted or re-flowed.  No further source-local context is needed: every cross-reference of every retained line resolves inside the two delivered spans.  Nothing was omitted from any retained span, so the package carries no omission record.  No retained line contains a citation, so the wrapper carries no bibliography and no bibliography entry.  No retained line contains a figure environment, an image-inclusion command or drawing code, so no image is delivered and the package declares no asset dependency.  Editorial formatting only, in addition: an AMS \texttt{amsart} preamble that restates the article's own declarations and macros used by the retained lines, namely the environments \texttt{lemma} on separate counters, together with the standard packages those lines need.  The retained environments print in this document as Lemma 1 in order of appearance, whereas the article prints the same items with the numbering of its own full document.  The complete native source of the article as distributed in the arXiv e-print for this exact version is bundled unchanged as \texttt{sources/210077/source0.tex}.
\begin{lemma}\label[lemma]{lemma appendix}
For every point $(w,v)$ away from the coordinate axes and every $h,k\in\R$,
\[
\mathfrak D_{\mathcal B}(w,v;h,k)
\leq-\frac12\mathsf A(w,v)k^2.
\]
\end{lemma}

\begin{proof}
When $p=2$, the assertion is immediate, so let
$1<p<2$. By evenness of $\mathcal B$ in the two variables, we may
assume that $w,v>0$. Set $A_0:=\mathsf A(w,v)<0$
and, for $(\xi,\eta)\in\R^2$, define
\[
\begin{aligned}
\Psi(\xi,\eta)
:=\mathcal B(w,v)
+\partial_w\mathcal B(w,v)(\xi-w)
+\partial_v\mathcal B(w,v)(\eta-v)
-\mathcal B(\xi,\eta)-\frac12A_0(\eta-v)^2.
\end{aligned}
\]
With $\xi=w+h$ and $\eta=v+k$, the desired inequality is equivalent to
$\Psi(\xi,\eta)\geq0$.

First, we may assume that both $\eta,\xi\geq0$. Indeed, by direct computation, since $w,v>0$,
\begin{align*}
\partial_w\mathcal B(w,v)
&=-\frac{C_pp}{p-1}(w+v)^{p-2}\bigl(w+(2-p)v\bigr)<0,\\
\partial_v\mathcal B(w,v)+A_0v
&=C_pp(p-2)v^2(w+v)^{p-3}\leq0.
\end{align*}
Consequently, for $\xi,\eta\geq0$,
\begin{align*}
\Psi(-\xi,\eta)-\Psi(\xi,\eta)
&=-2\xi\,\partial_w\mathcal B(w,v)\geq0,\\
\Psi(\xi,-\eta)-\Psi(\xi,\eta)
&=-2\eta\bigl(\partial_v\mathcal B(w,v)+A_0v\bigr)\geq0.
\end{align*}
Now, for fixed $\eta\geq0$, the function
$\xi\mapsto\Psi(\xi,\eta)$ is strictly convex on $[0,\infty)$, since
\[
\partial_{\xi\xi}\Psi(\xi,\eta)=-\mathsf A(\xi,\eta)>0.
\]
Define
\[
\Phi(\eta):=\min_{\xi\geq0}\Psi(\xi,\eta).
\]
At $\eta=v$, the minimum is attained at $\xi=w$, so
\[
\Phi(v)=\Phi'(v)=0.
\]

Suppose first that the minimizer $\xi=\xi(\eta)$ is positive. The
stationarity condition is
\[
\partial_w\mathcal B(\xi(\eta),\eta)=\partial_w\mathcal B(w,v).
\]
By setting
\[
s:=w+v,
\qquad S:=\xi+\eta,
\qquad \rho:=\frac vs,
\qquad R:=\frac\eta S,
\]
or equivalently
\[
(w,v)=s(1-\rho,\rho),
\qquad (\xi,\eta)=S(1-R,R),
\]
the stationarity condition becomes
\[
S^{p-1}\bigl(1-(p-1)R\bigr)
=s^{p-1}\bigl(1-(p-1)\rho\bigr).
\]
Solving for $S$ as a function of $R\in[0,1]$ gives
\[
S(R)
=s\left(\frac{1-(p-1)\rho}{1-(p-1)R}\right)^{1/(p-1)},
\qquad \xi(R):=(1-R)S(R),
\qquad H(R):=RS(R).
\]
Thus $(\xi(R),H(R))$ is the point on the stationary curve corresponding
to $R$.
Moreover, $H(\rho)=v$ and $H$ is strictly increasing, since
\[
H'(R)=S(R)\frac{1+(2-p)R}{1-(p-1)R}>0.
\]
The stationarity condition gives
\[
\Phi'(\eta)
=\partial_v\mathcal B(w,v)-A_0(\eta-v)-\partial_v\mathcal B(\xi,\eta),
\qquad
\Phi''(\eta)=-A_0-\Lambda_p(\xi,\eta).
\]
Along the stationary curve, set
\[
L(R):=\frac{\Lambda_p(\xi(R),H(R))}{-A_0}.
\]
Using the formulas for $\mathsf A$ and $\Lambda_p$, we obtain
\[
L(R)
=\left(\frac{1-(p-1)\rho}{1-(p-1)R}\right)^{\frac{p-2}{p-1}}
\frac1{\bigl(1+(2-p)R\bigr)\bigl(1+(2-p)\rho\bigr)}.
\]
In particular, $L(\rho)\leq1$, and one can prove that $L$ is decreasing.
Since $H$ is continuous and strictly increasing, it is a bijection from
$[0,1]$ onto $[0,\eta_b]$, where $\eta_b:=H(1)$.
For each $\zeta\in[0,\eta_b]$, let
\[
R_\zeta:=H^{-1}(\zeta).
\]
Equivalently, $R_\zeta$ is the unique parameter $R\in[0,1]$ such that
$H(R)=\zeta$. Thus
\[
H(R_\zeta)=\zeta,
\qquad R_{H(R)}=R.
\]
Since $H(\rho)=v$, we have in particular $R_v=\rho$ and $R_{\eta_b}=1$.
Setting $\widetilde L(\zeta):=L(R_\zeta)$, the identity for $\Phi''$ becomes
\[
\Phi''(\eta)=-A_0\bigl(1-\widetilde L(\eta)\bigr).
\]
We finally distinguish two cases.
If $v\leq\eta<\eta_b$, then $R_\eta\geq R_v=\rho$.
Since $L$ is decreasing, $\widetilde L(\eta)\leq L(\rho)\leq1$,
hence $\Phi''(\eta)\geq0$, and the identities at $\eta=v$ imply
$\Phi(\eta)\geq0$.

It remains to treat $0\leq\eta<v$. For $R\in[0,\rho]$, define
\[
G(R):=\int_{H(R)}^v\bigl(1-\widetilde L(t)\bigr)\,dt.
\]
Under the change of variables $t=H(u)$, the endpoints $t=H(R)$ and $t=v=H(\rho)$ correspond respectively to $u=R$ and $u=\rho$. Therefore
\[
G(R)=\int_R^\rho\bigl(1-L(u)\bigr)H'(u)\,du.
\]
Because $L$ is decreasing, $1-L$ is increasing and changes sign at most
once. Since
\[
G'(R)=-\bigl(1-L(R)\bigr)H'(R),
\]
$G$ has no interior minimum. Now $G(\rho)=0$, while
\[
L(R)H'(R)
=\frac{s\bigl(1-(p-1)\rho\bigr)}{1+(2-p)\rho}
\frac1{\bigl(1-(p-1)R\bigr)^2},
\]
and hence $G(0)\geq0$, thus $G(R)\geq0$ on $[0,\rho]$. Integrating the identity for $\Phi''$ twice and using $\Phi(v)=\Phi'(v)=0$ gives
\[
\Phi(\eta)
=-A_0\int_\eta^v(t-\eta)\bigl(1-\widetilde L(t)\bigr)\,dt.
\]
The two factors in the integrand are nondecreasing. Chebyshev's
inequality therefore yields
\begin{align*}
\int_\eta^v(t-\eta)\bigl(1-\widetilde L(t)\bigr)\,dt
&\geq\frac1{v-\eta}
\left(\int_\eta^v(t-\eta)\,dt\right)
\left(\int_\eta^v(1-\widetilde L(t))\,dt\right)\\
&=\frac{v-\eta}2G(R_\eta)\geq0.
\end{align*}
Since $-A_0>0$, this proves $\Phi(\eta)\geq0$ for $0\leq\eta<v$. Finally, at $\eta=\eta_b$ the interior minimizer reaches $\xi=0$.
By continuity,
\[
\Phi(\eta_b)\geq0,
\qquad \Phi'(\eta_b)\geq0.
\]
For $\eta\geq\eta_b$, the minimum is attained at $\xi=0$, and therefore
\[
\Phi(\eta)=\Psi(0,\eta),
\qquad \Phi''(\eta)=-A_0-\mathsf D(0,\eta).
\]
Direct calculation gives
\[
\mathsf D(0,\eta)=C_pp(p-1)\eta^{p-2},
\qquad
\Lambda_p(0,\eta)=\frac{C_pp}{3-p}\eta^{p-2}.
\]
Since $1<p<2$, we have
$\mathsf D(0,\eta)\leq\Lambda_p(0,\eta)$. Moreover,
$\Lambda_p(0,\eta)$ is decreasing in $\eta$ and
\[
\Lambda_p(0,\eta_b)=(-A_0)L(R_{\eta_b})=(-A_0)L(1)\leq-A_0.
\]
Consequently, for $\eta\geq\eta_b$,
\[
\mathsf D(0,\eta)
\leq\Lambda_p(0,\eta)
\leq\Lambda_p(0,\eta_b)
\leq-A_0,
\]
so $\Phi''(\eta)\geq0$. Together with the boundary values at $\eta_b$,
this yields $\Phi(\eta)\geq0$ for every $\eta\geq\eta_b$. This completes the proof.
\end{proof}

\end{document}
